% TOP_PROBLEM: {"id": 10100001, "title": "No percolation at the critical point on $\\mathbb{Z}^d$", "category_id": 19, "category": {"id": 19, "name": "probability", "display_name": "Probability", "description": "Problems involving probability theory, stochastic processes, and random structures.", "slug": "probability", "order_index": 19, "created_at": "2026-07-31T15:26:25.671Z"}, "status": "partially_solved"}
% FIRST_POSTED: null
% ENTRY_KIND: statement_only
% Research continuation compiled from the conversation, 27 September 2026.
% Root statement and all substantive progress preserved; final audit leaves the problem unresolved.
\documentclass[11pt]{article}
\usepackage[margin=25mm]{geometry}
\usepackage{fontspec}
\setmainfont{Latin Modern Roman}
\usepackage{amsmath,amssymb,mathtools}
\usepackage{unicode-math}
\setmathfont{Latin Modern Math}
\usepackage{xurl,hyperref}
\hypersetup{colorlinks=true,urlcolor=blue}
\setcounter{secnumdepth}{0}
\setlength{\parindent}{0pt}
\setlength{\parskip}{5pt}
\setlength{\emergencystretch}{3em}
\newcommand{\one}{\mathbf 1}
\begin{document}

\section*{\texorpdfstring{No percolation at the critical point on $\mathbb{Z}^d$}{No percolation at the critical point on \$\textbackslash{}mathbb\{Z\}\textasciicircum{}d\$}}\label{critical-percolation}\label{10100001}

\subsection{Definitions and mathematical statement}
Consider nearest-neighbour Bernoulli bond percolation on \(\mathbb Z^d\). Each edge is independently open with probability \(p\in[0,1]\). Let
\[
\theta(p)=\mathbb P_p(0\leftrightarrow\infty),
\qquad
p_c=\inf\{p:\theta(p)>0\}.
\]
The conjecture is
\[
\boxed{\theta(p_c)=0.}
\]
Equivalently, at criticality there is almost surely no infinite open cluster.

This is often called the \(\theta(p_c)=0\) conjecture, the critical percolation conjecture, the conjecture that there is no infinite cluster at criticality, or the continuity conjecture for the percolation probability at \(p_c\). There is no single universally fixed eponymous name.

The discussion below specializes to nearest-neighbour Bernoulli bond percolation on \(\mathbb Z^d\). A status checkpoint cited during the conversation was that the conjecture is known in \(d=2\), known in sufficiently high dimension (in particular \(d\ge 11\) in the cited 2026 source), and unresolved in dimensions \(3\le d\le10\); see \href{https://arxiv.org/html/2608.23661v1}{[S1]}.

\subsection{Short English statement}
At the exact threshold where percolation first becomes possible, is the probability that the origin belongs to an infinite open cluster still zero?

\subsection{Sources}
\begin{itemize}
\item[S1] Rapha\"el Cerf, 2026 preprint on critical percolation and finite-cluster tails.
\url{https://arxiv.org/html/2608.23661v1}.
\textit{Used in the conversation for the 2026 status checkpoint and the finite-cluster-tail alternative.}

\item[S2] Gady Kozma and Asaf Nitzan, 2024 preprint on critical percolation and finite-graph connectivity inequalities.
\url{https://arxiv.org/html/2401.12397v1}.
\textit{Used for local uniqueness, shell constructions, the conjectural gluing inequality, and the line/half-space route.}

\item[S3] Thomas M. Liggett, Roberto H. Schonmann, and Alan M. Stacey, ``Domination by product measures,'' \emph{Annals of Probability} 25 (1997).
\url{https://projecteuclid.org/journals/annals-of-probability/volume-25/issue-1/Domination-by-product-measures/10.1214/aop/1024404279.pdf}.
\textit{Used for finite-range dependent coarse-graining.}

\item[S4] Alexey Gladkov, 2024 preprint on decision-tree inequalities for percolation connectivity.
\url{https://arxiv.org/html/2408.08457v2}.
\textit{Used for the decision-tree mixing viewpoint and the adaptive-switching obstruction.}

\item[S5] Classical supercritical finite-cluster tail reference cited during the conversation.
\url{https://www.jstor.org/stable/2244302}.
\textit{Background for the distinction between strictly supercritical tail estimates and estimates needed from the weaker hypothesis \(\theta(p)>0\).}

\item[S6] Slab/percolation reference cited during the final audit.
\url{https://arxiv.org/pdf/1401.7130}.
\textit{Used to illustrate how negative information from an exposed path can obstruct full-space gluing and how quasi-planar geometry can help in slab settings.}
\end{itemize}

\end{document}
